\documentclass[a4paper,12pt]{article}
\usepackage[utf8]{inputenc}
\usepackage[T1]{fontenc}
\usepackage[polish]{babel}
\usepackage{geometry}
\usepackage{fancyhdr}
\usepackage{tcolorbox}
\usepackage{lmodern}
\usepackage{lastpage}
\usepackage{datetime}
\usepackage{titlesec}
\usepackage{changepage}
\usepackage[hidelinks]{hyperref}
\usepackage{graphicx}
\usepackage{listings}
\usepackage{url}
\usepackage{breakurl}
\usepackage{seqsplit}
\usepackage{adjustbox}
\usepackage{array}
\usepackage{longtable}
\usepackage{booktabs}
\usepackage{listingsutf8}
\lstset{inputencoding=utf8/latin2}

\geometry{margin=2.5cm}
\pagestyle{fancy}
\fancyhf{}
\fancyhead[L]{NFC SERVICES}
\fancyhead[R]{2025-07-08}
\fancyfoot[R]{\textbf{Strona \thepage\ z \pageref{LastPage}}}

% Konfiguracja dla długich linii kodu
\lstset{
    breaklines=true,
    breakatwhitespace=true,
    basicstyle=\footnotesize\ttfamily,
    frame=single,
    numbers=left,
    numberstyle=\tiny,
    stepnumber=1,
    numbersep=5pt,
    tabsize=2,
    showstringspaces=false,
    showtabs=false,
    captionpos=b,
    xleftmargin=2em,
    framexleftmargin=1.5em
}

\begin{document}
\vspace*{1cm}

\begin{center}
\begin{tcolorbox}[width=0.8\textwidth, sharp corners, boxrule=0.5pt]
    \centering
    \textbf{\Large Dokumentacja Techniczna \\ NFC Services} \\[1.5em]
    \textbf{Wersja:} 1.0 \\
    \textbf{Release:} A \\[1.5em]
    \textbf{Autor:} Michał Itrych \\[1.5em]
    \textbf{Data publikacji:} 08-07-2025
\end{tcolorbox}
\end{center}

\newpage
\tableofcontents
\newpage

\section{Wprowadzenie i przegląd systemu}

\subsection{Cel dokumentacji}
Niniejsza dokumentacja techniczna opisuje implementację systemu NFC Services dla aplikacji Flutter, umożliwiającego komunikację z kartami MIFARE DESFire i MIFARE Plus SL3.

\subsection{Architektura systemu}
System NFC Services składa się z następujących głównych komponentów:

\begin{itemize}
    \item \textbf{DesfireNfcServices} -- klasa obsługująca operacje na kartach MIFARE DESFire
    \item \textbf{PlusNfcServices} -- klasa obsługująca operacje na kartach MIFARE Plus SL3
    \item \textbf{Shared utilities} -- narzędzia do kryptografii, NFC oraz konwersji (ASCII, hex, bin).
\end{itemize}

\subsection{Wspierane typy kart}
System obsługuje następujące typy kart NFC:
\begin{itemize}
    \item \textbf{MIFARE DESFire} -- karty z zaawansowanymi funkcjami kryptograficznymi
    \item \textbf{MIFARE Plus SL3} -- karty z dodatkowymi mechanizmami bezpieczeństwa
\end{itemize}

\newpage

\subsection{Struktura katalogów}
\begin{lstlisting}[language=bash]
lib/
  nfc_services/
    desfire_services.dart
    plus_services.dart
    shared/
      crypto_utilities.dart
      nfc_utilities.dart
      binary_conversion.dart
\end{lstlisting}

\textbf{Opis plików:}
\begin{itemize}
    \item \texttt{desfire\_services.dart} -- Obsługa kart DESFire
    \item \texttt{plus\_services.dart} -- Obsługa kart MIFARE Plus
    \item \texttt{crypto\_utilities.dart} -- Narzędzia kryptograficzne
    \item \texttt{nfc\_utilities.dart} -- Podstawowe operacje NFC
    \item \texttt{binary\_conversion.dart} -- Konwersje binarne
    \item \texttt{desfire\_screen.dart} -- Przykład użycia DESFire
    \item \texttt{plus\_screen.dart} -- Przykład użycia MIFARE Plus
\end{itemize}

\subsection{Zależności zewnętrzne}
System korzysta z następujących bibliotek:

\begin{itemize}
    \item \texttt{flutter\_nfc\_kit} – do podstawowej komunikacji NFC,
    \item \texttt{convert} – do konwersji danych (np. między formatami hex, binarnym i ASCII),
    \item \texttt{encrypt} – do podstawowych operacji kryptograficznych,
    \item \texttt{pointycastle} – do zaawansowanych funkcji kryptograficznych.
\end{itemize}

Aby skorzystać z tych bibliotek, należy dodać je do pliku \texttt{pubspec.yaml}:

\begin{lstlisting}[language=bash]
dependencies:
  flutter_nfc_kit: ^3.3.1
  convert: ^3.0.1
  encrypt: ^5.0.1
  pointycastle: ^3.9.1
\end{lstlisting}


\subsection{Obsługiwane platformy}
System jest kompatybilny z:
\begin{itemize}
    \item \textbf{Android} -- SDK 26+ (Android 8.0+)
    \item \textbf{iOS} -- iOS 14.6+
\end{itemize}

\newpage

\section{Instalacja i konfiguracja}

\subsection{Wymagania systemowe}

\subsubsection{Wymagania sprzętowe}
\begin{itemize}
    \item \textbf{Urządzenie z NFC} -- telefon lub tablet z wbudowanym modułem NFC
    \item \textbf{Android} -- urządzenie z Android 8.0 (API 26) lub nowszym
    \item \textbf{iOS} -- iPhone z iOS 14.6 lub nowszym (iPhone 7 i nowsze)
\end{itemize}

\subsubsection{Wymagania programowe}
\begin{itemize}
    \item \textbf{Flutter SDK} -- wersja 3.7.2 lub nowsza
    \item \textbf{Dart SDK} -- wersja 3.0.0 lub nowsza
    \item \textbf{Android Studio} -- wersja 2022.1 lub nowsza (dla rozwoju Android)
    \item \textbf{Xcode} -- wersja 14.0 lub nowsza (dla rozwoju iOS)
    \item \textbf{Gradle} -- wersja 7.0 lub nowsza (dla Android)
\end{itemize}

\subsubsection{Wymagania dotyczące NFC}
\begin{itemize}
    \item \textbf{Standard NFC} -- urządzenie musi obsługiwać standard ISO/IEC 14443
    \item \textbf{Tryb czytnika} -- urządzenie musi obsługiwać tryb czytnika NFC
\end{itemize}

\subsection{Instalacja pakietu}

\subsubsection{Dodawanie zależności}
Aby zainstalować NFC Services w projekcie Flutter, dodaj następujące zależności do pliku \texttt{pubspec.yaml}:

\begin{lstlisting}
dependencies:
  flutter:
    sdk: flutter
  
  # NFC Services dependencies
  flutter_nfc_kit: ^3.3.1
  convert: ^3.0.1
  encrypt: ^5.0.1
  pointycastle: ^3.9.1
  
  # Additional dependencies for your app
  cupertino_icons: ^1.0.2
\end{lstlisting}

\subsubsection{Instalacja zależności}
Po dodaniu zależności, uruchom następujące komendy w terminalu:

\begin{lstlisting}[language=bash]
flutter pub get

flutter doctor

flutter run
\end{lstlisting}

\subsection{Konfiguracja uprawnień}

\subsubsection{Konfiguracja iOS}
Dodaj następujące klucze do pliku \texttt{ios/Runner/Info.plist}:

\begin{lstlisting}[language=xml]
<?xml version="1.0" encoding="UTF-8"?>
<!DOCTYPE plist PUBLIC "-//Apple//DTD PLIST 1.0//EN" "http://www.apple.com/DTDs/PropertyList-1.0.dtd">
<plist version="1.0">
<dict>
    <!-- Other keys... -->
    
    <key>NFCReaderUsageDescription</key>
	<string>NFC capabilities to function properly.</string>
</dict>
</plist>
\end{lstlisting}

Dodaj również następujące klucze do pliku \texttt{ios/Runner/Runner.entitlements}:

\begin{lstlisting}[language=xml]
<?xml version="1.0" encoding="UTF-8"?>
<!DOCTYPE plist PUBLIC "-//Apple//DTD PLIST 1.0//EN" "http://www.apple.com/DTDs/PropertyList-1.0.dtd">
<plist version="1.0">
<dict>
	<key>com.apple.developer.nfc.readersession.formats</key>
	<array>
		<string>TAG</string>
	</array>
</dict>
</plist>
\end{lstlisting}

\newpage

\subsection{Weryfikacja instalacji}

\subsubsection{Sprawdzenie dostępności NFC}
Aby sprawdzić czy NFC jest poprawnie skonfigurowane, użyj następującego kodu:

\begin{lstlisting}
import 'package:flutter_nfc_kit/flutter_nfc_kit.dart';

Future<void> checkNfcAvailability() async {
  try {
    final availability = await FlutterNfcKit.nfcAvailability;
    
    switch (availability) {
      case NFCAvailability.available:
        print('NFC is available and ready to use');
        break;
      case NFCAvailability.disabled:
        print('NFC is disabled - please enable in device settings');
        break;
      case NFCAvailability.not_supported:
        print('NFC is not supported on this device');
        break;
    }
  } catch (e) {
    print('Error checking NFC availability: $e');
  }
}
\end{lstlisting}

\newpage

\subsubsection{Test podstawowej funkcjonalności}
Utwórz prosty test aby sprawdzić czy wszystko działa:

\begin{lstlisting}
import 'package:flutter/material.dart';
import 'nfc_services/desfire_services.dart';

class NfcTestScreen extends StatefulWidget {
  @override
  _NfcTestScreenState createState() => _NfcTestScreenState();
}

class _NfcTestScreenState extends State<NfcTestScreen> {
  String _testResult = '';

  Future<void> _testNfcConnection() async {
    try {
      final desfire = DesfireNfcServices(debugMode: true);
      
      setState(() {
        _testResult = 'Testing NFC connection...';
      });
      
      // Test basic NFC availability
      final result = await desfire.readDesfire(
        'F40001', // Application ID
        '00',     // File ID
        '00',     // Key Number
        '00000000000000000000000000000000', // Auth Key
      );
      
      setState(() {
        if (result.isSuccess) {
          _testResult = 'NFC test successful! Data: ${result.data}';
        } else {
          _testResult = 'NFC test failed: ${result.errorCode?.description}';
        }
      });
    } catch (e) {
      setState(() {
        _testResult = 'Error during NFC test: $e';
      });
    }
  }

  @override
  Widget build(BuildContext context) {
    return Scaffold(
      appBar: AppBar(title: Text('NFC Test')),
      body: Padding(
        padding: EdgeInsets.all(16.0),
        child: Column(
          children: [
            ElevatedButton(
              onPressed: _testNfcConnection,
              child: Text('Test NFC Connection'),
            ),
            SizedBox(height: 16),
            Text(_testResult),
          ],
        ),
      ),
    );
  }
}
\end{lstlisting}

\subsection{Rozwiązywanie problemów z instalacją}

\subsubsection{Typowe problemy Android}
\begin{itemize}
    \item \textbf{Problem:} "NFC not available" na urządzeniu z NFC
    \begin{itemize}
        \item \textbf{Rozwiązanie:} Sprawdź czy NFC jest włączone w ustawieniach urządzenia
        \item \textbf{Rozwiązanie:} Upewnij się że aplikacja ma uprawnienia do NFC
    \end{itemize}
    
    \item \textbf{Problem:} "Gradle build failed"
    \begin{itemize}
        \item \textbf{Rozwiązanie:} Sprawdź czy wszystkie zależności są poprawnie dodane
        \item \textbf{Rozwiązanie:} Wykonaj \texttt{flutter clean} i \texttt{flutter pub get}
    \end{itemize}
\end{itemize}

\subsubsection{Typowe problemy iOS}
\begin{itemize}
    \item \textbf{Problem:} "NFC not supported" na iPhone
    \begin{itemize}
        \item \textbf{Rozwiązanie:} Upewnij się że używasz iPhone 7 lub nowszego
        \item \textbf{Rozwiązanie:} Sprawdź czy iOS jest w wersji 14.6+
    \end{itemize}
    
    \item \textbf{Problem:} "NFCReaderUsageDescription missing"
    \begin{itemize}
        \item \textbf{Rozwiązanie:} Dodaj opis użycia NFC do Info.plist
        \item \textbf{Rozwiązanie:} Upewnij się że klucz jest poprawnie sformatowany
    \end{itemize}
    
    \item \textbf{Problem:} "Entitlements error"
    \begin{itemize}
        \item \textbf{Rozwiązanie:} Sprawdź czy plik Runner.entitlements jest poprawnie skonfigurowany
        \item \textbf{Rozwiązanie:} Upewnij się że masz odpowiednie certyfikaty deweloperskie
    \end{itemize}
\end{itemize}

\subsubsection{Problem z zależnościami}
\begin{itemize}
    \item \textbf{Problem:} Konflikty wersji bibliotek
    \begin{itemize}
        \item \textbf{Rozwiązanie:} Użyj \texttt{flutter pub deps} aby sprawdzić zależności
        \item \textbf{Rozwiązanie:} Zaktualizuj Flutter do najnowszej wersji
    \end{itemize}
    
    \item \textbf{Problem:} "Package not found"
    \begin{itemize}
        \item \textbf{Rozwiązanie:} Sprawdź czy pubspec.yaml jest poprawnie sformatowany
        \item \textbf{Rozwiązanie:} Wykonaj \texttt{flutter pub cache repair}
    \end{itemize}
\end{itemize}

\section{Architektura i główne komponenty}

\subsection{Przegląd architektury}
System NFC Services ma modułową strukturę – jest podzielony na warstwę komunikacji z kartą, warstwę narzędzi kryptograficznych oraz warstwę prezentacji (np. przykładowe ekrany).

\subsection{Główne komponenty}

\begin{itemize}
    \item \textbf{DesfireNfcServices} – klasa odpowiedzialna za obsługę kart MIFARE DESFire, w tym autentykację, odczyt plików, obsługę błędów oraz logowanie.
    
    \item \textbf{PlusNfcServices} – klasa odpowiedzialna za obsługę kart MIFARE Plus SL3 (działa analogicznie do \texttt{DesfireNfcServices}, lecz oparta jest na innym protokole).
    
    \item \textbf{Narzędzia współdzielone (\texttt{shared/})} – zestaw narzędzi wykorzystywanych przez oba powyższe serwisy:
    \begin{itemize}
        \item \texttt{crypto\_utilities.dart} – funkcje kryptograficzne, m.in. szyfrowanie AES, generowanie liczb losowych, obliczanie CMAC oraz operacje na bajtach,
        \item \texttt{nfc\_utilities.dart} – funkcje wspomagające obsługę sesji NFC, takie jak sprawdzanie dostępności, wykrywanie kart, wysyłanie komend oraz obsługa błędów,
        \item \texttt{binary\_conversion.dart} – funkcje pomocnicze do konwersji danych binarnych i szesnastkowych (hex).
    \end{itemize}
    
\end{itemize}

\subsection{Schemat zależności}
\begin{itemize}
    \item \texttt{desfire\_services.dart} oraz \texttt{plus\_services.dart} korzystają z narzędzi z katalogu \texttt{shared/}.
    \item Wszystkie operacje NFC bazują na bibliotece \texttt{flutter\_nfc\_kit}.
\end{itemize}

\section{Moduł \texttt{DesfireNfcServices}}

\subsection{Opis klasy}
\texttt{DesfireNfcServices} to główna klasa odpowiedzialna za obsługę kart MIFARE DESFire. Realizuje pełny cykl komunikacji: od sprawdzenia dostępności NFC, przez wybór aplikacji i autentykację, aż po odczyt danych, obsługę błędów oraz generowanie logów.

\subsection{System obsługi wyników}
Każda metoda asynchroniczna zwraca wynik jako obiekt, który pochodzi z klasy \texttt{DesfireResult<T>}

\begin{itemize}
    \item \texttt{DesfireSuccess<T>} -- sukces operacji:
    \begin{itemize}
        \item \texttt{data} -- wynik operacji typu \texttt{T}
        \item \texttt{isSuccess} -- zawsze \texttt{true}
        \item \texttt{errorCode} -- zawsze \texttt{null}
        \item \texttt{message} -- zawsze \texttt{null}
        \item \texttt{exception} -- zawsze \texttt{null}
    \end{itemize}
    \item \texttt{DesfireError<T>} -- błąd operacji:
    \begin{itemize}
        \item \texttt{data} -- zawsze \texttt{null}
        \item \texttt{isSuccess} -- zawsze \texttt{false}
        \item \texttt{errorCode} -- kod błędu typu \texttt{DesfireErrorCode}
        \item \texttt{message} -- opcjonalny opis błędu (\texttt{String?})
        \item \texttt{exception} -- opcjonalny wyjątek (\texttt{Exception?})
    \end{itemize}
\end{itemize}

\textbf{Wspólne właściwości:}
\begin{itemize}
    \item \texttt{isSuccess} -- \texttt{bool} określający czy operacja zakończyła się sukcesem
    \item \texttt{data} -- wynik operacji lub \texttt{null} w przypadku błędu
    \item \texttt{errorCode} -- kod błędu lub \texttt{null} w przypadku sukcesu
    \item \texttt{message} -- opis błędu lub \texttt{null} w przypadku sukcesu
    \item \texttt{exception} -- wyjątek lub \texttt{null} w przypadku sukcesu
\end{itemize}

\newpage

\subsection{Typy wyników operacji}
System NFC Services wykorzystuje wzorzec \textbf{obsługi wyników} do spójnego zarządzania sukcesami i błędami operacji:

\subsubsection*{Struktura wyników}
\begin{lstlisting}[language=Java]
// Sukces
DesfireSuccess<String>(
  data: "48656C6C6F20576F726C642100000000"
)

// Blad
DesfireError<String>(
  error: DesfireErrorCode.authenticationFailed,
  message: "Authentication failed - check your key"
)
\end{lstlisting}

\subsubsection*{Sprawdzanie typu wyniku}
\begin{lstlisting}[language=Java]
final result = await desfire.readDesfire(...);

// Metoda 1:
if (result.isSuccess) {
    final data = result.data!; // Data is available
    print('Read data: $data');
} else {
    final error = result.errorCode!; // Error code is available
    print('Error: ${error.description}');
}

// Metoda 2:
switch (result) {
    case DesfireSuccess(data: final data):
        print('Read data: $data');
    case DesfireError(error: final error, message: final message):
        print('Error: ${error.description}');
        if (message != null) print('Details: $message');
}

// Metoda 3:
if (!result.isSuccess) {
    final error = result.errorCode!; // Error code is available
    print('Error: ${error.description}');
} else {
    final data = result.data!; // Data is available
    print('Read data: $data');
}
\end{lstlisting}

\newpage

\subsubsection*{Właściwości dostępne dla wszystkich wyników}
\begin{itemize}
    \item \texttt{isSuccess} -- \texttt{bool} - czy operacja się powiodła
    \item \texttt{data} -- \texttt{T?} - dane wyniku (null dla błędów)
    \item \texttt{errorCode} -- \texttt{DesfireErrorCode?} - kod błędu (null dla sukcesów)
\end{itemize}

\subsection{Konstruktor}
\begin{lstlisting}[language=Java]
DesfireNfcServices({
  bool debugMode = false, 
  int timeout = 7, 
  String alertMessage, 
  String multipleTagMessage
})
\end{lstlisting}

\textbf{Parametry:}
    \begin{itemize}
        \item \texttt{debugMode} (\texttt{bool}) -- włącza generowanie logów
        \item \texttt{timeout} (\texttt{int}) -- timeout operacji NFC (sekundy)
        \item \texttt{alertMessage} (\texttt{String}) -- komunikat podczas skanowania (tylko iOS)
        \item \texttt{multipleTagMessage} (\texttt{String}) -- komunikat przy wielu tagach (tylko iOS)
\end{itemize}

\subsection{Metody publiczne}
\begin{lstlisting}[language=Java]
Future<DesfireResult<String>> readDesfire(
  String applicationId, 
  String fileId, 
  String keyNumber, 
  String authKey
)
\end{lstlisting}

\textbf{Parametry:}
    \begin{itemize}
        \item \texttt{applicationId} (\texttt{String}) -- identyfikator aplikacji (6 znaków hex)
        \item \texttt{fileId} (\texttt{String}) -- identyfikator pliku (0-31)
        \item \texttt{keyNumber} (\texttt{String}) -- numer klucza (0-13)
        \item \texttt{authKey} (\texttt{String}) -- klucz autentykacji (32 znaki hex)
    \end{itemize}

\subsection{Typy pomocnicze}
\begin{itemize}
    \item \texttt{DesfireAuthResult} -- wynik autentykacji (kody, losowe liczby, klucz sesji)
    \item \texttt{DesfireFileSettings} -- ustawienia pliku (typ, tryb komunikacji, prawa dostępu)
    \item \texttt{DesfireFirstAuthResult} -- wynik pierwszego etapu autentykacji (RandomB)
\end{itemize}

\subsection{Metody prywatne}

\subsubsection*{Walidacja i obsługa błędów}
\begin{itemize}
    \item \texttt{DesfireErrorCode \_validateDesfireReadParameters(...)}\\
    Weryfikuje poprawność parametrów wejściowych (np. długość, format hex)
    
    \item \texttt{String \_getDesfireErrorMessage(String code)}\\
    Mapuje kod błędu na czytelny opis tekstowy
    
    \item \texttt{void \_validateCardResponse(...)}\\
    Sprawdza poprawność odpowiedzi karty (np. długość, status)
    
    \item \texttt{void \_log(String message, \{String level = 'INFO'\})}\\
    Loguje komunikaty (aktywne tylko w trybie \texttt{debugMode})
\end{itemize}

\subsubsection*{Operacje na karcie}
\begin{itemize}
    \item \texttt{Future<DesfireResult<void>> \_selectApplication(String applicationId)}\\
    Wysyła komendę wyboru aplikacji do karty
    
    \item \texttt{Future<DesfireResult<DesfireFileSettings>> \_getFileSettings(String fileId)}\\
    Pobiera ustawienia pliku (typ, prawa dostępu, tryb komunikacji)
    
    \item \texttt{Future<DesfireResult<String>> \_readEncryptedData(...)}\\
    Odczyt i deszyfracja danych z pliku
    
    \item \texttt{Future<String> \_decryptFileData(...)}\\
    Deszyfruje blok danych z pliku na podstawie klucza sesji i CMAC
\end{itemize}

\subsubsection*{Autentykacja}
\begin{itemize}
    \item \texttt{Future<DesfireResult<DesfireFirstAuthResult>> \_processFirstAuthStep(...)}\\
    Pierwszy etap autentykacji (pobranie i deszyfracja RandomB)
    
    \item \texttt{Future<DesfireResult<DesfireAuthResult>> \_processSecondAuthStep(...)}\\
    Drugi etap autentykacji (RandomA, sesja, weryfikacja)
    
    \item \texttt{Future<DesfireResult<DesfireAuthResult>> \_authenticateWithApplication(...)}\\
    Pełna autentykacja z obsługą retry
    
    \item \texttt{bool \_verifyRotatedRandomA(Uint8List received, Uint8List expected)}\\
    Porównuje wartości RandomA w sposób odporny na ataki czasowe
    
    \item \texttt{Uint8List \_generateSessionKey(Uint8List randomA, Uint8List randomB)}\\
    Generuje klucz sesji na podstawie RandomA i RandomB
    
    \item \texttt{String \_to3ByteLEHex(int value)}\\
    Zamienia liczbę na 3-bajtowy little-endian hex (do parametrów komend)
\end{itemize}

\newpage

\subsection{Obsługa błędów}
Każda metoda zwraca szczegółowy kod błędu (\texttt{DesfireErrorCode}) oraz opcjonalny opis i wyjątek.

\begin{table}[h!]
\centering
\begin{tabular}{|l|l|p{6cm}|}
\hline
\textbf{Kod błędu} & \textbf{Kod} & \textbf{Opis} \\
\hline
\texttt{success} & 00 & Operacja zakończona sukcesem \\
\hline
\texttt{invalidApplicationId} & DAID & Nieprawidłowy identyfikator aplikacji (musi być 6 znaków hex) \\
\hline
\texttt{invalidFileId} & DFID & Nieprawidłowy identyfikator pliku \\
\hline
\texttt{invalidKeyNumber} & DKN & Nieprawidłowy numer klucza \\
\hline
\texttt{invalidAuthKey} & DAK & Nieprawidłowy klucz autentykacji (musi być 32 znaki hex) \\
\hline
\texttt{timeoutError} & DT & Przekroczono limit czasu operacji \\
\hline
\texttt{nfcUnavailable} & DN & NFC nie jest dostępne \\
\hline
\texttt{sessionError} & DSN & Błąd sesji NFC \\
\hline
\texttt{unexpectedState} & DU & Nieoczekiwany stan aplikacji \\
\hline
\texttt{userCancelled} & DUC & Użytkownik anulował operację \\
\hline
\texttt{cardError} & CARD & Błąd karty \\
\hline
\texttt{authenticationFailed} & DAF & Błąd autentykacji \\
\hline
\texttt{sessionKeyError} & DSE & Błąd generowania klucza sesji \\
\hline
\texttt{cmacError} & DCM & Błąd obliczania CMAC \\
\hline
\texttt{dataDecryptionError} & DDD & Błąd deszyfracji danych \\
\hline
\texttt{invalidResponse} & DIR & Nieprawidłowa odpowiedź karty \\
\hline
\texttt{cardConnectionLost} & DCCL & Utracono połączenie z kartą - trzymaj kartę stabilnie \\
\hline
\texttt{communicationError} & DCE & Błąd komunikacji z kartą \\
\hline
\texttt{applicationNotFound} & DANF & Nie znaleziono aplikacji \\
\hline
\texttt{fileNotFound} & DFNF & Nie znaleziono pliku \\
\hline
\texttt{permissionDenied} & DPD & Odmowa dostępu \\
\hline
\texttt{conditionsNotSatisfied} & DCNS & Warunki użycia nie spełnione - kod polecenia niedostępny w obecnym stanie karty \\
\hline
\texttt{unknownError} & UE & Nieznany błąd. Spróbuj ponownie \\
\hline
\end{tabular}
\caption{Kody błędów \texttt{DesfireErrorCode} i ich znaczenie}
\end{table}

\clearpage

\newpage

\subsection{Przykład użycia}
\begin{lstlisting}[language=Java]
final desfire = DesfireNfcServices(debugMode: true);
final result = await desfire.readDesfire(
  'A1B2C3',
  '01',
  '00',
  '0123456789ABCDEF0123456789ABCDEF',
);

if (result.isSuccess) {
  print('Block data: ${result.data}');
  print('Operation successful: ${result.isSuccess}');
} else {
  print('Error code: ${result.errorCode}');
  print('Error message: ${result.message}');
  if (result.exception != null) {
    print('Exception: ${result.exception}');
  }
}
\end{lstlisting}

\subsection{Co może zwrócić \texttt{readDesfire}}

Metoda \texttt{readDesfire} zwraca obiekt typu \texttt{DesfireResult<String>}, który może być jednym z dwóch typów:

\subsubsection*{DesfireSuccess<String> - Sukces operacji}
\begin{itemize}
    \item \texttt{isSuccess} = \texttt{true}
    \item \texttt{data} = \texttt{String} - odczytane dane z pliku w formacie hex
    \item \texttt{errorCode} = \texttt{null}
    \item \texttt{message} = \texttt{null}
    \item \texttt{exception} = \texttt{null}
\end{itemize}

\textbf{Przykład sukcesu:}
\begin{lstlisting}[language=Java]
DesfireSuccess<String>(
  data: "48656C6C6F20576F726C642100000000" // "Hello World!" in hex
)
\end{lstlisting}

\newpage

\subsubsection*{DesfireError<String> - Błąd operacji}
\begin{itemize}
    \item \texttt{isSuccess} = \texttt{false}
    \item \texttt{data} = \texttt{null}
    \item \texttt{errorCode} = \texttt{DesfireErrorCode} - kod błędu
    \item \texttt{message} = \texttt{String?} - opcjonalny opis błędu
    \item \texttt{exception} = \texttt{Object?} - opcjonalny wyjątek
\end{itemize}

\textbf{Przykłady błędów:}
\begin{lstlisting}[language=Java]
// Blad: walidacja parametrow
DesfireError<String>(
  error: DesfireErrorCode.invalidApplicationId,
  message: "Invalid application ID (must be 6 hex characters)"
)

// Blad: NFC nie dostepne
DesfireError<String>(
  error: DesfireErrorCode.nfcUnavailable,
  message: "Please enable NFC in your device settings"
)

// Blad: autentykacja nie powiodla sie
DesfireError<String>(
  error: DesfireErrorCode.authenticationFailed,
  message: "Authentication failed - check your key"
)

// Blad : nieznany blad z wyjatkiem
DesfireError<String>(
  error: DesfireErrorCode.cardConnectionLost,
  message: "Card connection lost - please keep card steady"
)

// Blad: nieznany blad z wyjatkiem
DesfireError<String>(
  error: DesfireErrorCode.unknownError,
  message: "Unknown error. Try again",
  exception: TimeoutException("Operation timed out")
)
\end{lstlisting}

\newpage

\subsubsection*{Format zwracanych danych}
W przypadku sukcesu, \texttt{data} zawiera:
\begin{itemize}
    \item \textbf{Format:} String w formacie hex (32 znaki = 16 bajtów)
    \item \textbf{Przykład:} \texttt{"48656C6C6F20576F726C642100000000"}
    \item \textbf{Deszyfracja:} Dane są automatycznie deszyfrowane przed zwróceniem
    \item \textbf{Encoding:} Dane są w formacie little-endian zgodnie ze standardem DESFire
\end{itemize}

\subsubsection*{Obsługa błędów w aplikacji}
\begin{lstlisting}[language=Java]
final result = await desfire.readDesfire(appId, fileId, keyNum, authKey);

if (result.isSuccess) {
    // Sukces
    final blockData = result.data!;
    print('Read data: $blockData');
} else {
    // Blad
    final errorCode = result.errorCode!;
    final errorMessage = result.message;
    
    switch (errorCode) {
        case DesfireErrorCode.nfcUnavailable:
            showDialog('Enable NFC in device settings');
            break;
        case DesfireErrorCode.cardConnectionLost:
            showDialog('Keep the card steady near the phone');
            break;
        case DesfireErrorCode.authenticationFailed:
            showDialog('Invalid authentication key');
            break;
        default:
            showDialog('Error: ${errorMessage ?? errorCode.description}');
    }
}
\end{lstlisting}

\newpage

\section{Moduł \texttt{PlusNfcServices}}

\subsection{Opis klasy}
\texttt{PlusServices} to główna klasa odpowiedzialna za obsługę kart MIFARE Plus SL3. Klasa ta implementuje kompletny protokół uwierzytelniania i odczytu danych z kart MIFARE Plus, wykorzystując mechanizmy kryptograficzne AES128.


\subsection{System obsługi wyników}
Każda metoda asynchroniczna zwraca wynik jako obiekt, który pochodzi z klasy \texttt{PlusResult<T>}:

\begin{itemize}
    \item \texttt{PlusSuccess<T>} -- sukces operacji:
    \begin{itemize}
        \item \texttt{data} -- wynik operacji typu \texttt{T}
        \item \texttt{isSuccess} -- zawsze \texttt{true}
        \item \texttt{errorCode} -- zawsze \texttt{null}
        \item \texttt{message} -- zawsze \texttt{null}
        \item \texttt{exception} -- zawsze \texttt{null}
    \end{itemize}
    \item \texttt{PlusError<T>} -- błąd operacji:
    \begin{itemize}
        \item \texttt{data} -- zawsze \texttt{null}
        \item \texttt{isSuccess} -- zawsze \texttt{false}
        \item \texttt{errorCode} -- kod błędu typu \texttt{PlusErrorCode}
        \item \texttt{message} -- opcjonalny opis błędu (\texttt{String?})
        \item \texttt{exception} -- opcjonalny wyjątek (\texttt{Object?})
    \end{itemize}
\end{itemize}

\textbf{Wspólne właściwości:}
\begin{itemize}
    \item \texttt{isSuccess} -- \texttt{bool} określający czy operacja zakończyła się sukcesem
    \item \texttt{data} -- wynik operacji lub \texttt{null} w przypadku błędu
    \item \texttt{errorCode} -- kod błędu lub \texttt{null} w przypadku sukcesu
    \item \texttt{message} -- opis błędu lub \texttt{null} w przypadku sukcesu
    \item \texttt{exception} -- wyjątek lub \texttt{null} w przypadku sukcesu
\end{itemize}

\newpage

\subsection{Typy wyników operacji}
System NFC Services wykorzystuje wzorzec \textbf{obsługi wyników} do spójnego zarządzania sukcesami i błędami operacji:

\subsubsection*{Struktura wyników}
\begin{lstlisting}[language=Java]
// Sukces
PlusSuccess<String>(
  data: "48656C6C6F20576F726C642100000000"
)

// Blad
PlusError<String>(
  error: PlusErrorCode.authenticationFailed,
  message: "Authentication failed - check your key"
)
\end{lstlisting}

\subsubsection*{Sprawdzanie typu wyniku}
\begin{lstlisting}[language=Java]
final result = await plus.readPlus(...);

if (result.isSuccess) {
    final data = result.data!;
    print('Read data: $data');
} else {
    final error = result.errorCode!;
    print('Error: ${error.description}');
}

if (!result.isSuccess) {
    final error = result.errorCode!;
    print('Error: ${error.description}');
} else {
    final data = result.data!;
    print('Read data: $data');
}
\end{lstlisting}

\subsubsection*{Właściwości dostępne dla wszystkich wyników}
\begin{itemize}
    \item \texttt{isSuccess} -- \texttt{bool} - czy operacja się powiodła
    \item \texttt{data} -- \texttt{T?} - dane wyniku (null dla błędów)
    \item \texttt{errorCode} -- \texttt{PlusErrorCode?} - kod błędu (null dla sukcesów)
\end{itemize}

\newpage

\subsection{Konstruktor}
\begin{lstlisting}[language=Java]
PlusServices({
  bool debugMode = false, 
  int timeout = 10, 
  String alertMessage, 
  String multipleTagMessage
})
\end{lstlisting}

\textbf{Parametry:}
\begin{itemize}
    \item \texttt{debugMode} (\texttt{bool}) -- włącza szczegółowe logowanie
    \item \texttt{timeout} (\texttt{int}) -- timeout operacji NFC (sekundy)
    \item \texttt{alertMessage} (\texttt{String}) -- komunikat podczas skanowania (tylko iOS)
    \item \texttt{multipleTagMessage} (\texttt{String}) -- komunikat przy wielu tagach (tylko iOS)
\end{itemize}

\subsection{Metody publiczne}
\begin{lstlisting}[language=Java]
Future<PlusResult<String>> readPlus(
  int sector,
  int block,
  String keyType,
  String keyHex, {
  bool isPlaintext = true,
  bool requiresMacOnResponse = false,
})
\end{lstlisting}

\textbf{Parametry:}
\begin{itemize}
    \item \texttt{sector} (\texttt{int}) -- numer sektora (1-39)
    \item \texttt{block} (\texttt{int}) -- numer bloku (0-14)
    \item \texttt{keyType} (\texttt{String}) -- typ klucza ("Key A" lub "Key B")
    \item \texttt{keyHex} (\texttt{String}) -- klucz AES w formacie hex (32 znaki)
    \item \texttt{isPlaintext} (\texttt{bool}) -- czy odczytywać w trybie plaintext
    \item \texttt{requiresMacOnResponse} (\texttt{bool}) -- czy MAC jest wymagany w odpowiedzi
\end{itemize}

\subsection{Typy pomocnicze}
\begin{itemize}
    \item \texttt{PlusAuthResult} -- wynik autentykacji (kody, losowe liczby, transaction ID)
    \item \texttt{PlusSessionKeys} -- klucze sesji (kenc, kmac)
    \item \texttt{PlusReadResult} -- wynik odczytu bloku (kod, dane, opis)
    \item \texttt{KeyType} -- typ klucza (keyA, keyB)
\end{itemize}

\subsection{Metody prywatne}

\subsubsection*{Walidacja i obsługa błędów}
\begin{itemize}
    \item \texttt{PlusErrorCode \_validatePlusReadParameters(...)}\\
    Waliduje poprawność parametrów wejściowych (sektor, blok, typ klucza, format hex)
    
    \item \texttt{String \_getNfcErrorMessage(String code)}\\
    Mapuje kod błędu NFC na czytelny opis tekstowy
    
    \item \texttt{void \_log(String message, \{String level = 'INFO'\})}\\
    Loguje komunikaty (aktywne tylko w trybie \texttt{debugMode})
\end{itemize}

\subsubsection*{Autentykacja}
\begin{itemize}
    \item \texttt{Future<PlusResult<Uint8List>> \_processFirstAuthStep(...)}\\
    Pierwszy etap autentykacji (pobranie i deszyfracja RandomB)
    
    \item \texttt{Future<PlusResult<PlusAuthResult>> \_processSecondAuthStep(...)}\\
    Drugi etap autentykacji (RandomA, sesja, weryfikacja)
    
    \item \texttt{Future<PlusResult<PlusAuthResult>> \_authenticateSector(...)}\\
    Pełna autentykacja sektora z obsługą retry
    
    \item \texttt{PlusResult<PlusAuthResult> \_validateAuthResult(...)}\\
    Walidacja wyniku autentykacji i wyodrębnienie transaction ID
\end{itemize}

\subsubsection*{Generowanie kluczy sesji}
\begin{itemize}
    \item \texttt{PlusResult<PlusSessionKeys> \_generateSessionKeys(...)}\\
    Generuje klucze sesji (kenc, kmac) na podstawie danych autentykacji
\end{itemize}

\subsubsection*{Operacje na blokach}
\begin{itemize}
    \item \texttt{Future<PlusResult<PlusReadResult>> \_readBlockData(...)}\\
    Odczytuje dane z bloku z obsługą CMAC
    
    \item \texttt{String \_decryptBlockData(...)}\\
    Deszyfruje dane bloku przy użyciu kluczy sesji
    
    \item \texttt{String \_extractBlockData(...)}\\
    Wyodrębnia dane konkretnego bloku z pełnej odpowiedzi
\end{itemize}

\subsubsection*{Generowanie komend}
\begin{itemize}
    \item \texttt{String \_generateAuthCommand(...)}\\
    Generuje komendę autentykacji (70h)
    
    \item \texttt{String \_generateBlockReadCommandCode(...)}\\
    Generuje kod komendy odczytu bloku (31h, 32h, 33h)
    
    \item \texttt{String \_generateBlockReadCommand(...)}\\
    Generuje pełną komendę odczytu bloku
    
    \item \texttt{String \_calculateCmac(...)}\\
    Oblicza CMAC dla autentykacji komendy
\end{itemize}

\subsection{Obsługa błędów}
Każda metoda zwraca szczegółowy kod błędu (\texttt{PlusErrorCode}) oraz opcjonalny opis i wyjątek. Kody błędów są mapowane na czytelne komunikaty, co ułatwia diagnostykę i obsługę wyjątków w aplikacji.

\begin{table}[h!]
\centering
\begin{tabular}{|l|l|p{6cm}|}
\hline
\textbf{Kod błędu} & \textbf{Kod} & \textbf{Opis} \\
\hline
\texttt{success} & 00 & Operacja zakończona sukcesem \\
\hline
\texttt{invalidSector} & PPS & Nieprawidłowy numer sektora (musi być 1-39) \\
\hline
\texttt{invalidBlock} & PPB & Nieprawidłowy numer bloku (musi być 0-14) \\
\hline
\texttt{invalidKeyType} & PPKT & Nieprawidłowy typ klucza \\
\hline
\texttt{invalidKeyHex} & PPK & Nieprawidłowy format klucza AES hex (musi być 32 znaki hex) \\
\hline
\texttt{timeoutError} & PPT & Przekroczono limit czasu operacji \\
\hline
\texttt{nfcUnavailable} & PPN & NFC nie jest dostępne \\
\hline
\texttt{sessionError} & PPSN & Błąd sesji NFC \\
\hline
\texttt{unexpectedState} & PPU & Nieoczekiwany stan aplikacji \\
\hline
\texttt{userCancelled} & PPUC & Użytkownik anulował operację \\
\hline
\texttt{cardError} & CARD & Błąd karty \\
\hline
\texttt{authenticationFailed} & PPAF & Błąd autentykacji \\
\hline
\texttt{sessionKeyError} & PPSE & Błąd generowania klucza sesji \\
\hline
\texttt{cmacError} & PPCM & Błąd obliczania CMAC \\
\hline
\texttt{dataDecryptionError} & PPDD & Błąd deszyfracji danych \\
\hline
\texttt{invalidResponse} & PPIR & Nieprawidłowa odpowiedź karty \\
\hline
\texttt{cardConnectionLost} & PCCL & Utracono połączenie z kartą - trzymaj kartę stabilnie \\
\hline
\texttt{communicationError} & PCE & Błąd komunikacji z kartą \\
\hline
\texttt{unknownError} & UE & Nieznany błąd. Spróbuj ponownie \\
\hline
\end{tabular}
\caption{Kody błędów \texttt{PlusErrorCode} i ich znaczenie}
\end{table}

\newpage

\subsection{Przykład użycia}
\begin{lstlisting}[language=Java]
final plus = PlusServices(debugMode: true);
final result = await plus.readPlus(
  1,
  0,
  "Key A",
  "0123456789ABCDEF0123456789ABCDEF",
  isPlaintext: true,
  requiresMacOnResponse: false,
);

if (result.isSuccess) {
  print('Block data: ${result.data}');
  print('Operation successful: ${result.isSuccess}');
} else {
  print('Error code: ${result.errorCode}');
  print('Error message: ${result.message}');
  if (result.exception != null) {
    print('Exception: ${result.exception}');
  }
}
\end{lstlisting}

\subsection{Co może zwrócić \texttt{readPlus}}

Metoda \texttt{readPlus} zwraca obiekt typu \texttt{PlusResult<String>}, który może być jednym z dwóch typów:

\subsubsection*{PlusSuccess<String> - Sukces operacji}
\begin{itemize}
    \item \texttt{isSuccess} = \texttt{true}
    \item \texttt{data} = \texttt{String} - odczytane dane z bloku w formacie hex
    \item \texttt{errorCode} = \texttt{null}
    \item \texttt{message} = \texttt{null}
    \item \texttt{exception} = \texttt{null}
\end{itemize}

\textbf{Przykład sukcesu:}
\begin{lstlisting}[language=Java]
PlusSuccess<String>(
  data: "48656C6C6F20576F726C642100000000" // "Hello World!" in hex
)
\end{lstlisting}

\newpage

\subsubsection*{PlusError<String> - Błąd operacji}
\begin{itemize}
    \item \texttt{isSuccess} = \texttt{false}
    \item \texttt{data} = \texttt{null}
    \item \texttt{errorCode} = \texttt{PlusErrorCode} - kod błędu
    \item \texttt{message} = \texttt{String?} - opcjonalny opis błędu
    \item \texttt{exception} = \texttt{Object?} - opcjonalny wyjątek
\end{itemize}

\textbf{Przykłady błędów:}
\begin{lstlisting}[language=Java]
// Blad: walidacja parametrow
PlusError<String>(
  error: PlusErrorCode.invalidSector,
  message: "Invalid sector (must be 1-39)"
)

// Blad: NFC nie dostepne
PlusError<String>(
  error: PlusErrorCode.nfcUnavailable,
  message: "Please enable NFC in your device settings"
)

// Blad: autentykacja nie powiodla sie
PlusError<String>(
  error: PlusErrorCode.authenticationFailed,
  message: "Authentication failed - check your key"
)

// Blad: utracono polaczenie z karta
PlusError<String>(
  error: PlusErrorCode.cardConnectionLost,
  message: "Card connection lost - please keep card steady"
)

// Blad: nieznany blad z wyjatkiem
PlusError<String>(
  error: PlusErrorCode.unknownError,
  message: "Unknown error. Try again",
  exception: TimeoutException("Operation timed out")
)
\end{lstlisting}

\newpage

\subsubsection*{Format zwracanych danych}
W przypadku sukcesu, \texttt{data} zawiera:
\begin{itemize}
    \item \textbf{Format:} String w formacie hex (32 znaki = 16 bajtów)
    \item \textbf{Przykład:} \texttt{"48656C6C6F20576F726C642100000000"}
    \item \textbf{Deszyfracja:} Dane są automatycznie deszyfrowane przed zwróceniem
\end{itemize}

\subsubsection*{Obsługa błędów w aplikacji}
\begin{lstlisting}[language=Java]
final result = await plus.readPlus(sector, block, keyType, keyHex);

if (result.isSuccess) {
    // Sukces - uzyj result.data
    final blockData = result.data!;
    print('Read data: $blockData');
} else {
    // Blad - sprawdz kod i wiadomosc
    final errorCode = result.errorCode!;
    final errorMessage = result.message;
    
    switch (errorCode) {
        case PlusErrorCode.nfcUnavailable:
            showDialog('Enable NFC in device settings');
            break;
        case PlusErrorCode.cardConnectionLost:
            showDialog('Keep the card steady near the phone');
            break;
        case PlusErrorCode.authenticationFailed:
            showDialog('Invalid authentication key');
            break;
        default:
            showDialog('Error: ${errorMessage ?? errorCode.description}');
    }
}
\end{lstlisting}

\newpage

\section{Moduły współdzielone}

System NFC Services zawiera moduły współdzielone, które zapewniają wspólne funkcjonalności dla obu serwisów (MIFARE DESFire i MIFARE Plus). Te moduły zawierają narzędzia kryptograficzne, konwersje binarne oraz podstawowe operacje NFC.

\subsection{Moduł \texttt{binary\_conversion.dart}}

\subsubsection{Opis klasy}
\texttt{BinaryConversion} to klasa narzędziowa do konwersji danych między różnymi formatami. Zapewnia metody do konwersji ciągów hex na formaty binarne, hex i ASCII z obsługą wyboru zakresu bajtów oraz odwracania bajtów.

\subsubsection{Główne funkcjonalności}
\begin{itemize}
    \item \textbf{Konwersja do formatu binarnego} -- \texttt{toBinary()} konwertuje hex na ciąg binarny z odstępami między bajtami
    \item \textbf{Konwersja do formatu hex} -- \texttt{toHex()} zwraca hex z odstępami między bajtami w formacie wielkich liter
    \item \textbf{Konwersja do ASCII} -- \texttt{toASCII()} konwertuje hex na czytelny tekst ASCII z zastąpieniem nieprawidłowych znaków kropkami
\end{itemize}

\subsubsection{Przykłady użycia}
\begin{lstlisting}[language=Java]
// Konwersja do binarnego
final binary = BinaryConversion.toBinary("12345678", 0, 1);
// Wynik: "00010010 00110100"

// Konwersja do hex
final hex = BinaryConversion.toHex("12345678", 0, 1);
// Wynik: "12 34"

// Konwersja do ASCII
final ascii = BinaryConversion.toASCII("48656C6C6F", 0, 4);
// Wynik: "Hello"

\end{lstlisting}

\subsubsection{Obsługa błędów}
Klasa zawiera szczegółową walidację parametrów wejściowych:
\begin{itemize}
    \item Sprawdzenie czy ciąg nie jest pusty
    \item Weryfikacja parzystej długości ciągu hex
    \item Walidacja poprawności znaków hex
    \item Kontrola zakresu pozycji bajtów
\end{itemize}

\subsection{Moduł \texttt{crypto\_utilities.dart}}

\subsubsection{Opis klasy}
\texttt{CryptoUtilities} to klasa narzędziowa dla operacji kryptograficznych. Zapewnia metody do szyfrowania AES, generowania liczb losowych, manipulacji bajtami oraz obliczania CMAC dla operacji MIFARE Plus.

\subsubsection{Główne funkcjonalności}

\paragraph{Szyfrowanie AES}
\begin{itemize}
    \item \texttt{aesEncrypt()} -- szyfruje dane hex z kluczem i IV w trybie AES-CBC
    \item \texttt{aesDecrypt()} -- deszyfruje dane hex z kluczem i IV w trybie AES-CBC
    \item Obsługa trybu CBC bez paddingu dla operacji NFC
    \item Walidacja parametrów (32 znaki hex dla klucza i IV)
\end{itemize}

\paragraph{Generowanie liczb losowych}
\begin{itemize}
    \item \texttt{generateSecureRandom()} -- generuje kryptograficznie bezpieczne losowe bajty
    \item Używa \texttt{Random.secure()} dla operacji kryptograficznych
    \item Typowo używane do generowania RandomA i RandomB w protokołach NFC
\end{itemize}

\paragraph{Manipulacja bajtami}
\begin{itemize}
    \item \texttt{rotateLeft()} -- obraca tablicę bajtów o jedną pozycję w lewo
    \item \texttt{rotateRight()} -- obraca tablicę bajtów o jedną pozycję w prawo
    \item \texttt{xorBytes()} -- wykonuje operację XOR na dwóch tablicach bajtów
    \item \texttt{padToBlockSize()} -- dodaje padding PKCS\#7 do danych
    \item \texttt{removePadding()} -- usuwa padding PKCS\#7 z danych
\end{itemize}

\paragraph{Konwersje liczbowe}
\begin{itemize}
    \item \texttt{intToLittleEndianHex()} -- konwertuje liczbę na 4-znakowy little-endian hex
    \item \texttt{littleEndianHexToInt()} -- konwertuje little-endian hex na liczbę
    \item Używane do formatowania komend NFC gdzie wymagany jest little-endian
\end{itemize}

\paragraph{Obliczanie CMAC}
\begin{itemize}
    \item \texttt{computeAesCmac()} -- oblicza AES-CMAC z określonym rozmiarem wyjścia
    \item Obsługuje rozmiary 128, 192 i 256 bitów
    \item Używane do weryfikacji integralności w operacjach MIFARE Plus
\end{itemize}

\subsubsection{Przykłady użycia}
\begin{lstlisting}[language=Java]
// Szyfrowanie AES
final encrypted = CryptoUtilities.aesEncrypt(
    "48656C6C6F20576F726C64", // "Hello World" in hex
    "0123456789ABCDEF0123456789ABCDEF", // 32-char key
    "00000000000000000000000000000000"  // 32-char IV
);

// Generowanie losowych bajtow
final randomBytes = CryptoUtilities.generateSecureRandom(16);

// Obliczanie CMAC
final cmac = CryptoUtilities.computeAesCmac(
    "0123456789ABCDEF0123456789ABCDEF", // key
    "48656C6C6F20576F726C64", // data
    128 // output bits
);

// Konwersja little-endian
final leHex = CryptoUtilities.intToLittleEndianHex(258);
// Wynik: "0201" (nie "0102")

// XOR bajtow
final result = CryptoUtilities.xorBytes(
    Uint8List.fromList([1, 2, 3]),
    Uint8List.fromList([4, 5, 6])
);
// Wynik: [5, 7, 5]
\end{lstlisting}

\subsubsection{Obsługa błędów}
Klasa zawiera własny typ wyjątku \texttt{CryptoException} oraz szczegółową walidację:
\begin{itemize}
    \item Sprawdzenie poprawności parametrów AES (długość, format hex)
    \item Walidacja parametrów CMAC (rozmiar wyjścia, format danych)
    \item Kontrola rozmiaru tablic bajtów dla operacji XOR
    \item Weryfikacja paddingu PKCS\#7
\end{itemize}

\subsubsection{Zależności zewnętrzne}
Moduł wykorzystuje następujące biblioteki:
\begin{itemize}
    \item \texttt{convert} -- konwersje hex/binary
    \item \texttt{encrypt} -- operacje kryptograficzne AES
    \item \texttt{pointycastle} -- obliczanie CMAC i zaawansowane funkcje kryptograficzne
\end{itemize}

\subsection{Moduł \texttt{nfc\_utilities.dart}}

\subsubsection{Opis klasy}
\texttt{nfc\_utilities.dart} to moduł zawierający podstawowe narzędzia do obsługi operacji NFC. Zapewnia funkcje do sprawdzania dostępności NFC, wykrywania kart, wysyłania komend oraz obsługi błędów komunikacji z kartami.

\subsubsection{Główne funkcjonalności}
\begin{itemize}
    \item \textbf{Sprawdzanie dostępności NFC} -- \texttt{isNfcAvailable()} weryfikuje czy NFC jest dostępne, wyłączone lub nieobsługiwane
    \item \textbf{Wykrywanie kart} -- \texttt{pollForCard()} skanuje i wykrywa karty NFC z obsługą timeout'ów
    \item \textbf{Wysyłanie komend} -- \texttt{sendCommand()} i \texttt{sendCommands()} wysyłają komendy hex do karty
    \item \textbf{Walidacja odpowiedzi} -- \texttt{isSuccessfulResponse()} sprawdza poprawność odpowiedzi karty
    \item \textbf{Obsługa błędów} -- \texttt{NfcErrorCode} i \texttt{NfcResult} zapewniają spójną obsługę błędów
    \item \textbf{Debug logging} -- \texttt{debugPrint()} i \texttt{debugPrintWithLevel()} dla diagnostyki
\end{itemize}

\subsubsection{Obsługiwane typy kart}
Moduł obsługuje następujące typy kart NFC:
\begin{itemize}
    \item \texttt{mifarePlus} -- karty MIFARE Plus SL3
    \item \texttt{mifareDesfire} -- karty MIFARE DESFire
\end{itemize}

\subsubsection{Przykład użycia}
\begin{lstlisting}[language=Java]
// Sprawdzanie dostepnosci NFC
final nfcResult = await isNfcAvailable();
if (!nfcResult.isSuccess) {
    print('NFC error: ${nfcResult.message}');
    return;
}

// Wykrywanie karty
final cardResult = await pollForCard(
    debugMode: true,
    timeout: 10,
    alertMessage: "Please place card near device",
    multipleTagMessage: "Multiple cards detected"
);
}
\end{lstlisting}

\end{document}